\documentclass[11pt,a4paper]{article}

\usepackage[T1]{fontenc}
\usepackage[utf8]{inputenc}
\usepackage{lmodern}
\usepackage[margin=2.35cm]{geometry}
\usepackage{amsmath,amssymb}
\usepackage{array,booktabs,longtable,tabularx}
\usepackage{graphicx}
\usepackage{microtype}
\usepackage{placeins}
\usepackage{natbib}
\usepackage[hidelinks]{hyperref}

\graphicspath{{figures/}}
\setlength{\parindent}{0pt}
\setlength{\parskip}{0.55em}
\setlength{\emergencystretch}{2em}
\setcounter{tocdepth}{2}
\newcommand{\unit}[1]{\,\mathrm{#1}}

\title{QBO composite and metrics methodology}
\author{}
\date{\today}

\begin{document}

\maketitle

\section{Overview}

This document defines a composite framework for characterising the
Quasi-Biennial Oscillation (QBO).  Westerly and easterly onsets are identified
from the detrended, deseasonalised and smoothed 30-hPa zonal wind averaged over
$5^\circ$S--$5^\circ$N.  Atmospheric fields composited about these onsets are
used to diagnose latitudinal width, period, amplitude, phase, vertical extent
and descent rate.  Individual cycles are compared with the composite using a
coherence and correlation metrics.

This methodology is a variation of the QBOi methodology, using 30 hPa as the reference pressure level and
extending the compositing method to include latitudinal dependence. See \citet{pahlavan2021, simpson2025,
yu2017, lee2024} for other examples of a QBO composite methodology.

\begin{table}[htbp]
\centering
\caption{Summary of the diagnostic framework.}
\begin{tabularx}{\textwidth}{@{}
>{\raggedright\arraybackslash}p{0.24\textwidth}X@{}}
\toprule
Diagnostic & Methodological definition \\
\midrule
Reference dates & Linearly interpolated zero crossings of the centred
five-month running mean of deseasonalised and linearly detrended 30-hPa zonal
wind averaged over $5^{\circ}$S--$5^{\circ}$N.  The default direction is
negative to positive (westerly onset); crossings separated by less than five
months are averaged. \\
Composite & Mean deseasonalised anomaly over complete $\pm15$-month cycle
windows. Users may exclude specified onset years and apply linear detrending. \\
Latitudinal width & Full width at half maximum of a fitted parabolic-cylinder
structure at each pressure level. \\
Period, amplitude and phase & Sinusoidal fits to the composite lag profiles at each pressure level.
The default QBO period is the mean period from the easterly- and
westerly-onset zonal wind composites at 30~hPa. \\
Basic amplitude & Half the maximum-minus-minimum composite range over lag for
each variable and location. \\
Vertical extent & Pressure level at 10\%-of-maximum of the fitted amplitude profile (with interpolation). \\
Descent rate & Mean downward displacement of the onset-connected zero-wind
contour between successive lags, with a pressure-resolved profile calculated
from its crossing lag at each pressure. \\
Cycle-to-cycle coherence & Squared lag correlation between the raw data and its composite over the default
$\pm15$-month cycle window at each pressure level.\\
Composite variance explained & Fraction of cycle and lag variance explained
by the composite at each pressure level.\\
\bottomrule
\end{tabularx}
\end{table}


\begin{table}
\centering
\caption{Key symbols. Detailed definitions in the text.}
\begin{tabularx}{\textwidth}{@{}
>{\raggedright\arraybackslash}p{0.24\textwidth}X@{}}
\toprule
Symbol & Definition \\
\midrule
$u(t,p,\phi)$ & Monthly-mean zonal-mean zonal wind. \\
$u_{30}(t)$ & $5^{\circ}$S--$5^{\circ}$N mean wind at 30~hPa. \\
$r(t)$ &  Detrended, deseasonalised and smoothed reference series derived from $u_{30}$. \\
$x$, $x'_c$ & An atmospheric field and its deseasonalised anomaly in QBO
cycle $c$. \\
$C_x$ & Composite mean of $x'_c$; for example, $C_u$ is the zonal-wind composite. \\
$\widehat{C}_x(\phi,p)$ & Fitted latitudinal profile of $C_x$ at a specified lag. \\
$\widehat{C}_x(\tau,p)$ & Fitted lag profile of $C_x$ at a specified latitude. \\
\bottomrule
\end{tabularx}
\end{table}

\section{Composite methodology}
\label{sec:composite}

\subsection{Reference series and onset dates}

Let $u(t,p,\phi)$ be monthly zonal-mean zonal wind, where $t$ is time, $p$ is pressure and $\phi$ is latitude.
The reference series is formed from the 30-hPa wind averaged over $5^{\circ}$S--$5^{\circ}$N,
\begin{equation}
u_{30}(t)=\left\langle u(t,30\unit{hPa},\phi)\right\rangle_{\phi},
\qquad -5^{\circ}\leq\phi\leq5^{\circ},
\end{equation}
where angle brackets denote the latitude average.  Wind is interpolated to
30~hPa if that pressure is not available directly. All interpolations are linear in log-pressure coordinates
(where applicable).

For monthly analysis, the calendar-month climatology and a linear trend are removed from $u_{30}$ before
applying a centred five-month running mean.  The result is the reference timeseries $r(t)$, from which
reference dates are defined by identifying onset of westerly or easterly winds.

To calculate the QBO period with sub-monthly precision, there is an option to estimate the period using daily
means, a cyclic 31-day-smoothed calendar-day climatology and a centred 152-day mean.  The 29~February
climatology is the mean of those for 28~February and 1~March. A linear trend is removed before smoothing, as
for the monthly calculation.

A westerly onset is a negative-to-positive crossing of $r(t)$; an easterly onset is positive to negative.  In
either case the onset time is interpolated between the two bracketing months where a zero-crossing occurs. If
two crossings occur within 5 (nominal) months of each other, the midpoint is used as the reference date. This
prevents a short reversal from defining an additional onset. Figure~\ref{fig:reference} shows the
westerly-onset (top) and easterly-onset (bottom) dates and the reference timeseries calculated from ERA5.
All ERA5 composite figures exclude cycles with onset years 2015 and 2020, corresponding to the 2016 and 2020 QBO
disruptions. The disruptions are included in the cycle-to-cycle variability plots.

\begin{figure}[htbp]
\centering
\includegraphics[width=0.96\textwidth]{reference_timeseries_exclude_2015_2020.pdf}
\caption{Monthly 30-hPa reference timeseries $r(t)$ for westerly and easterly onsets from ERA5.
Dashed lines mark zero crossings. A merged pair of nearby crossings would be shown by
an orange interval and dotted original crossings; none occur in this dataset.}
\label{fig:reference}
\end{figure}

\subsection{Cycle alignment and compositing}

For the composites, we recommend using monthly-mean data. Each interpolated onset date is aligned to its nearest
monthly sample and a 31-sample window ($\tau=-15,\ldots,15$ months) is extracted.  QBO cycles without a
complete window are omitted.  Given a field $x$, the calendar-month climatology is removed at each pressure
and latitude to form the cycle anomaly $x'_c(\tau, p, \phi)$. Cycle fields may be linearly detrended but by
default are not. QBO cycles can be excluded by their onset year (e.g. to exclude disruptions). For cycle anomaly
$x'_c$, the composite is
\begin{equation}
C_x(\tau,p,\phi)=\frac{1}{N_c}\sum_{c=1}^{N_c}
x'_c(\tau,p,\phi),
\end{equation}
where $N_c$ is the number of QBO cycles. The onset direction is left implicit
in $C_x$; westerly- and easterly-onset composites are formed separately.
Figure~\ref{fig:lag-height} shows the westerly-onset zonal wind and
temperature composites.

\begin{figure}[htbp]
\centering
\includegraphics[width=0.90\textwidth]{lag_height_composite_exclude_2015_2020.pdf}
\caption{Equatorial westerly-onset composites from ERA5. Filled contours are temperature
anomalies (K); line contours are zonal-wind anomalies
($\mathrm{m\,s^{-1}}$). The grey dashed line marks the 30-hPa reference
pressure.}
\label{fig:lag-height}
\end{figure}

\FloatBarrier

\section{QBO metrics}
\label{sec:metrics}

\subsection{Latitudinal width}

Latitudinal structure is diagnosed from the composite profile at zero lag.  At each
pressure level we fit a sum of parabolic-cylinder functions
\citep{paracyl},
\begin{equation}
\widehat{C}_x(\phi,p)=\sum_j a_j(p)
D_{n_j}\!\left(\frac{\phi-c(p)}{s(p)}\right),
\end{equation}
where $a_j(p)$ is a component amplitude, $c(p)$ and $s(p)$ are the common
centre and latitude scale, and $j$ indexes the fitted components. Fits use $30^{\circ}$S--$30^{\circ}$N by default,
with $|c(p)|\leq5^{\circ}$ and $2^{\circ}\leq s(p)\leq30^{\circ}$.

For zonal wind, we use the Gaussian $D_0$ by default. For temperature, we use
the sum of $D_0$ and $D_2$. The temperature fit should also be used for ozone
($\mathrm{O}_3$). The function $D_2$ looks like a central Gaussian with two side-lobes, with $D_0$ included to fit the central amplitude.
This allows a better fit with both the equatorial and subtropical contributions.

The latitudinal width is calculated as the full width at half maximum (FWHM) of the absolute fitted equatorial
lobe. For completeness, we also report the scale $s(p)$ (which, if equivalent to the FWHM, would be half its
value).  A fit is accepted when its RMS residual, normalised by the full-cycle lag--latitude RMS after removal
of the lag mean, is at most 0.5 and its log-FWHM differs by at most 0.1 from interpolation between adjacent
presssure levels. This is designed to reject poor fits between the QBO jets and in the upper troposphere/lower
stratosphere where the QBO amplitude weakens. Both thresholds are configurable. Figure~\ref{fig:width-fits}
shows the latitudinal width calculation for zonal wind at 50 hPa and temperature at 30 hPa from ERA5.

\begin{figure}[htbp]
\centering
\includegraphics[width=\textwidth]{latitudinal_width_fits_exclude_2015_2020.pdf}
\caption{Parabolic-cylinder fits to zonal wind near 50~hPa and temperature
near 30~hPa from ERA5.  Solid and dashed markers show FWHM and fitted scale,
respectively.}
\label{fig:width-fits}
\end{figure}

\subsection{Period, amplitude, phase, and vertical extent}

At each pressure level, the composite lag profile is fitted with
\begin{equation}
\widehat{C}_x(\tau,p)=b(p)+A(p)\sin\!\left[
\frac{2\pi\tau}{P}+\theta(p)\right],
\end{equation}
where $b(p)$ is the offset, $A(p)>0$ is the amplitude, $P$ is the QBO period
and $\theta(p)\in[-\pi,\pi)$ is the phase.

A basic half-range amplitude is also calculated directly from each
composite variable:
\begin{equation}
A_x^{\mathrm{range}}(p,\phi)=\frac{1}{2}\left[
\max_\tau C_x(\tau,p,\phi)-\min_\tau C_x(\tau,p,\phi)\right].
\end{equation}

The QBO period is the mean period obtained from fits to the easterly- and westerly-onset zonal wind composites
$C_u$ at 30 hPa, interpolated in pressure if necessary. This period is held fixed at all pressure levels for
subsequent calculations of the amplitude and phase for a given variable. The resulting amplitude and phase
profiles therefore describe one common oscillation frequency. The QBO period is continuous rather than rounded
to the sampling period of the underlying composite. If only monthly raw data is supplied, the sub-monthly
precision of the period is therefore inferred. 

The smallest and largest individual QBO periods provide a simple uncertainty
range around the sine-fit period.  Each individual period is the interval
between consecutive same-direction zero crossings at 30~hPa.  Both onset
directions contribute to the reported range.  Excluding an onset year removes
the cycle that starts at that crossing from the range.

If daily data is supplied, the QBO period can optionally be calculated by applying the same compositing and
fitting procedure to daily raw data.  Nominal-month parameters are converted using 365.25/12 days per month:
the reference smoothing is 152 days and the $\pm15$-month cycle window becomes $\pm457$ days.  Each
interpolated onset time is aligned to its nearest daily sample.  The unsmoothed, deseasonalised 30-hPa wind
at the equator is then composited and fitted as a function of lag in days. As
for monthly cycle fields, these daily anomalies are not detrended by default.
The final reported period is an average of the calculation for a
westerly- and easterly-onset composite. Figure~\ref{fig:daily-period} shows the fit calculation on
daily-sampled composites.

\begin{figure}[htbp]
\centering
\includegraphics[width=\textwidth]{daily_period_fits_exclude_2015_2020.pdf}
\caption{Daily ERA5 30-hPa composites and sine fits.  The periods for different choices of onset
differ by 53.5 days and average to 853.9 days (28.05 months).  Individual
zero-crossing periods range from 18.47 to 37.46 months. The calculation
excludes the 2015 and 2020 onset years.}
\label{fig:daily-period}
\end{figure}

The vertical extent of the QBO is defined as the linearly interpolated pressure level where the amplitude
$A(p)$ is 10\% of its maximum:
\begin{equation}
	A(p_{\mathrm{top}})=0.1\max_p\{A(p)\}.
\end{equation}
No value is assigned if it the threshold is not met in the range of pressure levels provided.
Figure~\ref{fig:phase-amplitude} shows the phase and both amplitude profiles for zonal wind and temperature using
the daily-sampled QBO period of 28.05 months with the vertical extent and the reference height (30 hPa)
indicated.

\begin{figure}[htbp]
\centering
\includegraphics[width=0.94\textwidth]{phase_amplitude_exclude_2015_2020.pdf}
\caption{Amplitude and phase profiles for a QBO period of 28.05 months from ERA5.
Blue solid lines show sine-fit amplitude. Purple dashed lines show half-range amplitude.
Grey dashed and green dash--dot lines mark 30~hPa and the vertical extent, respectively.}
\label{fig:phase-amplitude}
\end{figure}

\FloatBarrier

\subsection{Descent rate}

The mean descent rate follows \citet{Schenzinger2017},
adapted to pressure coordinates and the zonal-wind composite $C_u$.  For each
onset direction, the zero-wind contour passing through the 30-hPa onset is
followed by linear interpolation.  Its pressure $p_0(\tau)$ gives the
downward displacement between successive monthly lags,
\begin{equation}
v_{d,i}=\frac{p_0(\tau_{i+1})-p_0(\tau_i)}
{\tau_{i+1}-\tau_i},
\end{equation}
where positive $v_{d,i}$ indicates downward motion.  The reported mean is the
arithmetic mean of its positive monthly values.

A pressure-resolved rate is obtained from the lag $\tau_0(p)$ at which the
same contour crosses each pressure level.  At interior levels a centred
difference gives
\begin{equation}
v_d(p_j)=\frac{p_{j+1}-p_{j-1}}
{\tau_0(p_{j+1})-\tau_0(p_{j-1})},
\end{equation}
with one-sided differences at the ends.  Only the monotonically descending
part of the contour is retained, with no extrapolation beyond its diagnosed
pressure range.  The calculation is applied separately to easterly- and
westerly-onset composites between 1 and 200~hPa by default.  The positive
monthly rates from each individual QBO cycle are interpolated onto the common
pressure levels to calculate the cycle-to-cycle standard deviation.

\begin{figure}[htbp]
\centering
\includegraphics[width=0.62\textwidth]{descent_rate_exclude_2015_2020.pdf}
\caption{Pressure-resolved downward rates of the onset-connected zero-wind
contours in the ERA5 easterly- and westerly-onset composites.  Legend values
are the mean monthly displacements.  Shading spans one standard deviation
across individual QBO cycles and the grey dashed line marks 30~hPa.}
\label{fig:descent}
\end{figure}

\subsection{Cycle-to-cycle coherence}

For each equatorial zonal-wind cycle $u'_c$ (with a $\pm15$-month window by default) and pressure level,
coherence is the squared Pearson correlation with $C_u$ across lag,
\begin{equation}
\mathcal{C}_{c,p}=\operatorname{corr}_{\tau}
\left[u'_c(\tau,p),C_u(\tau,p)\right]^2.
\end{equation}
The coherence is bounded by $0 \le \mathcal{C}_{c, p} \le 1$ and measures similarity of temporal evolution
between individual cycles and the QBO composite. The cycle mean is weighted by pressure-layer thickness in log
pressure,
\begin{equation}
w_p=\frac{|\nabla\log p|}{\sum_p|\nabla\log p|}.
\end{equation}
The vertically averaged coherence of cycle $c$ is therefore
\begin{equation}
\overline{\mathcal{C}}_c=\sum_p w_p\mathcal{C}_{c,p}.
\end{equation}
The composite variance explained at pressure $p$ is
\begin{equation}
V_p=1-\frac{\operatorname{var}_{c,\tau}
[u'_c(\tau,p)-C_u(\tau,p)]}
{\operatorname{var}_{c,\tau}[u'_c(\tau,p)]}.
\end{equation}
Figure~\ref{fig:coherence} shows both metrics with the QBO wind anomalies.
These diagnostics span 0.5--200~hPa by default.

\begin{figure}[p]
\centering
\includegraphics[width=\textwidth]{cycle_coherence.pdf}
\caption{QBO-cycle anomalies, pressure-resolved coherence and composite
variance explained. The variance profile repeats beside both pressure panels.
The calculation includes the 2015 and 2020 onset years.
Vertical dashed lines mark westerly onsets; the
horizontal grey dashed line marks 30~hPa. Green dash--dot and dotted lines in
the upper panel mark the top and bottom vertical extent, respectively.}
\label{fig:coherence}
\end{figure}

Alongside coherence, we use two metrics to capture the cycle-to-cycle variability in QBO amplitude, structure
and timing. For any field $q(\tau,p)$, define the weighted mean
\begin{equation}
\left\langle q\right\rangle_W
=\frac{1}{N_\tau}\sum_\tau\sum_p w_p q(\tau,p),
\end{equation}
and let $\mu_c=\langle u'_c\rangle_W$ and $\mu_{C_u}=\langle C_u\rangle_W$.  The
whole-pattern correlation is
\begin{equation}
R_c=\frac{\left\langle
(u'_c-\mu_c)(C_u-\mu_{C_u})\right\rangle_W}
{\sqrt{\left\langle(u'_c-\mu_c)^2\right\rangle_W
       \left\langle(C_u-\mu_{C_u})^2\right\rangle_W}}.
\end{equation}
It is a signed measure on $[-1,1]$: positive values indicate similar
lag--pressure structure and timing, while negative values indicate an
opposite pattern.

The RMS-amplitude ratio is
\begin{equation}
A_c^{\mathrm{RMS}}
=\frac{\sqrt{\left\langle (u'_c)^2\right\rangle_W}}
       {\sqrt{\left\langle C_u^2\right\rangle_W}}.
\end{equation}
A value of one gives equal RMS amplitude; values above and below one indicate
stronger and weaker cycles, respectively.  Both metrics are dimensionless.
Figure~\ref{fig:cycle-summary} plots $R_c$ against $A_c^{\mathrm{RMS}}$ for
each cycle, with onset year shown by colour and the point label.
At each lag and pressure, values outside the standard 1.5-IQR limits are
excluded from the cycle standard deviation; the composite itself retains all
cycles.

\begin{figure}[htbp]
\centering
\includegraphics[width=\textwidth]{cycle_summary.pdf}
\caption{Westerly-onset QBO cycles at 30~hPa and their amplitude, timing and structure metrics from ERA5.
The left panel shows the individual onset-aligned cycles (coloured by onset year), their composite (black),
and the composite plus or minus two standard deviations across cycles (grey shading).  The right panel shows
whole-pattern correlation against RMS-amplitude ratio; colour and labels give onset year.  Dashed lines mark
zero lag in the left panel and unit correlation and amplitude ratio in the right panel. The standard deviation
excludes pointwise values outside the standard 1.5-IQR limits.}
\label{fig:cycle-summary}
\end{figure}


\FloatBarrier

\section{Configurable parameters}
\label{sec:parameters}

\small
\begin{longtable}{@{}
>{\raggedright\arraybackslash}p{0.29\textwidth}
>{\raggedright\arraybackslash}p{0.19\textwidth}
>{\raggedright\arraybackslash}p{0.44\textwidth}@{}}
\caption{User-configurable methodology parameters and defaults.}
\label{tab:parameters}\\
\toprule
Parameter & Default & Role \\
\midrule
\endfirsthead
\multicolumn{3}{l}{\tablename\ \thetable\ continued}\\
\toprule
Parameter & Default & Role \\
\midrule
\endhead
\midrule
\multicolumn{3}{r}{Continued on next page}\\
\endfoot
\bottomrule
\endlastfoot
\multicolumn{3}{@{}l}{\textit{Reference timeseries}}\\[0.25em]
Reference pressure & 30 hPa & Onset-defining level; interpolated if absent. \\
Reference latitude band & $5^{\circ}$S--$5^{\circ}$N & Band averaged for the
reference timeseries. \\
Reference smoothing & 5 months & Centred running-mean for monthly data. Daily data uses an equivalent
152-sample running mean. \\
Daily climatology smoothing & 31 days & Cyclic mean of the calendar-day
climatology. \\
Crossing direction & Westerly & Negative-to-positive; easterly is the reverse. \\
Crossing merge tolerance & 5 months & Minimum separation between onset events. \\
\addlinespace
\multicolumn{3}{@{}l}{\textit{Composites}}\\[0.25em]
Cycle half-window & 15 months & 31 monthly lags or $\pm457$ daily lags. \\
Composite deseasonalisation & On & Remove the calendar-month climatology. \\
Composite detrending & Off & Remove each field's linear trend before cycle
extraction. \\
Excluded onset years & None & Omit specified QBO cycles from the composite. \\
\addlinespace
\multicolumn{3}{@{}l}{\textit{Period}}\\[0.25em]
Daily period latitude & $0^{\circ}$ & Latitude to fit for the daily period estimate. \\
Initial period guess & 28 months & Initial value; the fitted period is bounded
between half and twice this value. \\
\addlinespace
\multicolumn{3}{@{}l}{\textit{Latitudinal width}}\\[0.25em]
Parabolic-cylinder orders & $D_0$ (u); $D_0+D_2$ ($T$, $O_3$) &
Variable-dependent default. \\
Fit latitude range & $30^{\circ}$S--$30^{\circ}$N & Calculation range. \\
Fit pressure range & 1--200 hPa & Calculation range. \\
Maximum normalised fit error & 0.5 & Residual-to-full-cycle-RMS threshold; 0.3
is used for wind in the ERA5 illustration. \\
Maximum log-width discontinuity & 0.1 & Vertical-continuity threshold for
log-FWHM. \\
\addlinespace
\multicolumn{3}{@{}l}{\textit{Phase \& amplitude}}\\[0.25em]
Supplied QBO period & None & If absent, estimate from the 30-hPa composites;
the daily period estimate may instead be supplied. \\
Period reference pressure & 30 hPa & Height of the QBO period fit. \\
Phase/amplitude pressure range & 1--200 hPa & Calculation range. \\
\addlinespace
\multicolumn{3}{@{}l}{\textit{Vertical extent}}\\[0.25em]
Vertical-extent threshold & 0.10 & Fraction of maximum amplitude at the upper
boundary. \\
\addlinespace
\multicolumn{3}{@{}l}{\textit{Descent rate}}\\[0.25em]
Descent pressure range & 1--200 hPa & Calculation range. \\
Descent reference pressure & 30 hPa & Selects the onset-connected zero-wind contour. \\
\addlinespace
\multicolumn{3}{@{}l}{\textit{Coherence}}\\[0.25em]
Coherence pressure range & 0.5--200 hPa & Calculation range. \\
\end{longtable}
\normalsize

\clearpage

\section{Illustration with ERA5}

The figures use ERA5 for 1979--2022: daily wind for the optional period
estimate and monthly data otherwise.  The composites exclude cycles with
onsets in 2015 and 2020, representing the 2016 and 2020 disruptions.
The period diagnostics use the same exclusions.  The cycle-to-cycle
variability diagnostics retain both onset years.
After these exclusions, 15 westerly and 17 easterly onsets have complete cycle windows.  Table
\ref{tab:example-results} illustrates the resulting metrics; the values are
not prescribed observational targets.

\begin{table}[htbp]
\centering
\caption{ERA5 results. Composite and period metrics exclude the disruption-onset
years; variability metrics include them.}
\label{tab:example-results}
\begin{tabularx}{0.94\textwidth}{@{}Xr@{}}
\toprule
Quantity & ERA5 illustration \\
\midrule
Complete westerly/easterly cycle windows & 15 / 17 \\
Complete westerly variability windows & 17 \\
Monthly-sampled QBO period (easterly-/westerly-onset) & 28.54 / 27.34 months \\
Daily-sampled QBO period (easterly-/westerly-onset) & 28.93 / 27.18 months \\
Daily-sampled QBO period (average of both onsets) & 28.05 months / 853.9 days \\
Individual zero-crossing period range & 18.47--37.46 months \\
Easterly--westerly period difference & 53.5 days \\
Top vertical extent, easterly/westerly composites & 0.83 / 0.79 hPa \\
Bottom vertical extent, easterly/westerly composites & 89.44 / 87.18 hPa \\
Mean descent rate, easterly-/westerly-onset composite & 4.87 / 7.31 $\mathrm{hPa\,month^{-1}}$ \\
Mean westerly-onset cycle coherence & 0.52 \\
Westerly-onset wind FWHM near 50 hPa & $22.8^{\circ}$ \\
Westerly-onset temperature FWHM near 30 hPa & $16.6^{\circ}$ \\
\bottomrule
\end{tabularx}
\end{table}

\begin{thebibliography}{9}
\bibitem[Schenzinger et al.(2017)]{Schenzinger2017}
V. Schenzinger, S. Osprey, L. Gray, and N. Butchart,
``Defining metrics of the Quasi-Biennial Oscillation in global climate
models,'' \emph{Geoscientific Model Development}, 10, 2157--2168, 2017.
\href{https://doi.org/10.5194/gmd-10-2157-2017}
{doi:10.5194/gmd-10-2157-2017}.

\bibitem[Pahlavan et al.(2021)]{pahlavan2021}
H. A. Pahlavan, Q. Fu, J. M. Wallace, and G. N. Kiladis,
``Revisiting the Quasi-Biennial Oscillation as seen in ERA5. Part I:
Description and momentum budget,'' \emph{Journal of the Atmospheric
Sciences}, 78, 673--691, 2021.
\href{https://doi.org/10.1175/JAS-D-20-0248.1}
{doi:10.1175/JAS-D-20-0248.1}.

\bibitem[Simpson et al.(2025)]{simpson2025}
I. R. Simpson et al.,
``The path toward vertical grid options for the Community Atmosphere Model
version 7: The impact of vertical resolution on the QBO and tropical waves,''
\emph{Journal of Advances in Modeling Earth Systems}, 17, e2025MS004957,
2025.
\href{https://doi.org/10.1029/2025MS004957}
{doi:10.1029/2025MS004957}.

\bibitem[Yu et al.(2017)]{yu2017}
C. Yu, X. Xue, J. Wu, T. Chen, and H. Li,
``Sensitivity of the quasi-biennial oscillation simulated in WACCM to the
phase speed spectrum and the settings in an inertial gravity wave
parameterization,'' \emph{Journal of Advances in Modeling Earth Systems}, 9,
389--403, 2017.
\href{https://doi.org/10.1002/2016MS000824}
{doi:10.1002/2016MS000824}.

\bibitem[Lee et al.(2024)]{lee2024}
H.-K. Lee, H.-Y. Chun, J. H. Richter, I. R. Simpson, and R. R. Garcia,
``Contributions of parameterized gravity waves and resolved equatorial waves
to the QBO period in a future climate of CESM2,'' \emph{Journal of
Geophysical Research: Atmospheres}, 129, e2024JD040744, 2024.
\href{https://doi.org/10.1029/2024JD040744}
{doi:10.1029/2024JD040744}.

\bibitem[Weisstein(n.d.)]{paracyl}
E. W. Weisstein,
``Parabolic Cylinder Function,'' \emph{MathWorld---A Wolfram Resource}.
\href{https://mathworld.wolfram.com/ParabolicCylinderFunction.html}
{mathworld.wolfram.com/ParabolicCylinderFunction.html} (accessed 17 July
2026).
\end{thebibliography}

\end{document}
